\styledchapter{Measure Spaces}
\section{Introduction}


\begin{example}[Non-measurable event]
	Consider a uniform probability measure (measure of $1$) on $[0,2\pi)$. For example, \[
		\P\left[[0,1]\right] = \frac{1}{2\pi}, \quad \P\left[[0,1] \cup [2,3]\right] = \frac{1}{\pi}
	.\] 

	Define the function $T$ by \[
		T(x) = \begin{cases}
			x+1, \quad x+1 < 2\pi \\
			x+1 - 2\pi, \quad x+1 \ge 2\pi
		\end{cases}
	.\] 
	Thinking of a circle, the operator moves the point by one unit. With this operator, we can define an equivalence relation: we will say $x \sim y$ iff $T(x) = T(y)$ or $T(y) = T(x)$.

	The equivalence relation leads to a partition $(A_\lambda:\lambda)$; i.e., \[
		A_\lambda \cap A_\lambda' = \O, \ \lambda \ne \lambda', \quad \cup_\lambda A_\lambda = [0,2\pi)
	\] such that $x,y \in A_\lambda \iff x \sim y$.

	For every $\lambda$, select an element $x_\lambda \in A_\lambda$ and collect them in $X = \{x_\lambda : \lambda\}$\boxednote{$X$ can be called a representative set.}.\boxednote{Here we use the Axiom of Choice.} Define the translation \[
		X+i = \{x_\lambda + i : \lambda\}
	.\] 

	Note that $\{X+i \operatorname{mod} 2\pi: i \in \Z\}$ is a partition of $[0,2\pi)$. 

	We know that
	\begin{align*}
		1 = \P\left[[0,2\pi)\right] &= \sum\limits_{i \in \Z}^{} \P\left[(X+i)\right]  \\
									&= \sum\limits_{i \in \Z}^{} \P\left[X\right]
	.\end{align*}
	Then what is $\P\left[X\right]$? It is not well-defined; $X$ is not a measurable event.
\end{example}

We need to restrict ourselves to measurable subsets to have our notion of probability be well-defined; hence measure theory is necessary.

\begin{definition}[$\sigma$-algebra]
	Let $\Omega$ denote the sample space and $2^{\Omega} := \{\text{subsets of } \Omega\}$.

	Given $\Omega$, $\mathscr{F} \subseteq 2^{\Omega}$ is a $\sigma$-field ($\sigma$-algebra) if
	\begin{enumerate}[label=(\roman*)]
		\item $\O \in \mathscr{F}$
		\item If $A \in \mathscr{F}$, then $A^{c} \in \mathscr{F}$
		\item If $A_n \in \mathscr{F}$, then $\cup_n A_n \in \mathscr{F}$ (countable union).
	\end{enumerate}
\end{definition}

Trivial examples include $2^{\Omega}$ and $\{\O,\Omega\}$. 

\begin{definition}[Generated $\sigma$-field]
	Given $\mathscr{E} \subseteq 2^{\Omega}$, define $\sigma(\mathscr{E})$ to be the smallest $\sigma$-field containing $\mathscr{E}$. In other words, for any $\sigma$-field $\mathscr{F} \supseteq \mathscr{E}$, then $\sigma(\mathscr{E}) \subseteq \mathscr{F}$.

	Note that $\sigma(\mathscr{E})$ is well-defined, because we can put \[
		\sigma(\mathscr{E}) = \bigcap \{\mathscr{F} \subseteq 2^{\Omega}: \mathscr{F} \text{ is a } \sigma\text{-field and } \mathscr{F} \supseteq \mathscr{E}\}
	.\] We use the following lemma, whose proof is immediate.
\end{definition}

\begin{lemma}
	If $\mathscr{F}_\lambda$ is a $\sigma$-field for each $\lambda$, then $\cap_\lambda \mathscr{F}_\lambda$ is a $\sigma$-field.
\end{lemma}

\begin{definition}[Monotone sequences of sets and their limits]
	We say that $A_n$ is an increasing sequence ($A_n \uparrow$) if \[
		A_1 \subseteq A_2 \subseteq \ldots
	\] and similarly for decreasing sequences of sets.

	If $A_n \uparrow$, define $\lim\limits_{n \to \infty} A_n := \cup_n A_n$, and if $A_n \downarrow$, define $\lim\limits_{n \to \infty} A_n := \cap_n A_n$.
\end{definition}

We can define the limit of an arbitrary, non-monotone sequence.

\begin{definition}[$\liminf$ and $\limsup$ of sequences]
	For arbitrary $A_n$, define \[
		\liminf\limits_{n \to \infty} A_n = \cup_k \cap_{n \ge k} A_n, \quad \limsup\limits_{n \to \infty} A_n = \cap_k \cup_{n \ge k} A_n
	.\] If $\liminf\limits_{n \to \infty} A_n = \limsup\limits_{n \to \infty} A_n$, then the limit exists and is the common value.
\end{definition}

\begin{remark}
	Thinking hard, these definitions make sense. Alternatively, 
	\begin{align*}
		x \in \liminf\limits_{n \to \infty} A_n &\iff \exists\ k \text{ s.t. } \forall\ n \ge k ,\  x\in A_n \\
												&\iff x \text{ belongs to all but finitely many } A_n \\
		x \in \limsup_{n \to \infty} A_n &\iff \forall\ k, \exists\ n \ge k \text{ s.t. } x \in A_n \\
										 &\iff x \text{ belongs to infinitely many } A_n
	.\end{align*}
	Furthermore, $\liminf\limits_{n \to \infty} A_n \subseteq \limsup\limits_{n \to \infty} A_n$.
\end{remark}

\begin{definition}[$\pi$-system, monotone class, $\lambda$-system, field]
	We say that $\mathscr{P}$ is a $\pi$-system if it is closed under intersection; i.e., for any $A,B \in \mathscr{P}$, $A \cap B \in \mathscr{P}$.

	We say that $\mathscr{M}$ is a monotone class if for any monotone sequence $(A_n)$, $\lim\limits_{n \to \infty} A_n \in \mathscr{M}$. Monotone classes are closed under the limit operation.

	We say that $\mathscr{L}$ is a $\lambda$-system if 
	\begin{enumerate}[label=(\roman*)]
		\item $\Omega \in \mathscr{L} $,
		\item If $A, B \in \mathscr{L}$ such that $A \supseteq B$, then $A \setminus B \in \mathscr{L}$, and 
		\item If $(A_n)$ is increasing sequence in $\mathscr{L}$, then $\lim\limits_{n \to \infty} A_n \in \mathscr{L}$.\boxednote{This assumption is weaker than the one in the definition of a $\sigma$-field.} $\mathscr{L}$ is closed under countable unions, but only for increasing sequences.
	\end{enumerate}

	$\mathscr{A}$ is a field (algebra) if 
	\begin{enumerate}[label=(\roman*)]
		\item $\O \in \mathscr{A}$,
		\item $\forall\  A \in \mathscr{A}$, $A^{c}\in \mathscr{A}$, and 
		\item $\forall\  A, B \in \mathscr{A}$, $A \cup B \in \mathscr{A}$.
	\end{enumerate}
\end{definition}

The first and fourth definitions deal with finite operations; the second and third include limits.

\begin{theorem}
	(Monotone class.)\boxednote{$m(\mathscr{A})$ is defined to be the smallest monotone class containing $\mathscr{A}$.} If $\mathscr{A}$ is a field, then $m(\mathscr{A}) = \sigma(\mathscr{A})$.
\end{theorem}

\begin{proof}
	The inclusion $m(\mathscr{A}) \subseteq \sigma(\mathscr{A})$ is straightforward.
	By definition, $\sigma(\mathscr{A})$ contains $\mathscr{A}$.
	Also, $\sigma(\mathscr{A})$ is a monotone class.
	Since $m(\mathscr{A})$ is the smallest monotone class containing $\mathscr{A}$, we get
	$m(\mathscr{A}) \subseteq \sigma(\mathscr{A})$.

	For the reverse inclusion, it suffices to show that $m(\mathscr{A})$ is a $\sigma$-field
	containing $\mathscr{A}$. It contains $\mathscr{A}$ by definition.

	(i) Since $\mathscr{A}$ is a field, $\O \in \mathscr{A} \subseteq m(\mathscr{A})$.

	(ii) (Closure under complements.) Define
	\[
		M_0 := \{A \in m(\mathscr{A}) : A^{c} \in m(\mathscr{A})\}.
	\]
	We claim $M_0$ is a monotone class containing $\mathscr{A}$.

	\emph{(a) $\mathscr{A} \subseteq M_0$.}
	If $A \in \mathscr{A}$, then $A^c \in \mathscr{A} \subseteq m(\mathscr{A})$, hence $A \in M_0$.

	\emph{(b) $M_0$ is a monotone class.}
	Suppose $A_n \uparrow A$ with each $A_n \in M_0$.
	Then $A_n^c \downarrow A^c$, and since each $A_n^c \in m(\mathscr{A})$ and $m(\mathscr{A})$
	is a monotone class, we have $A^c \in m(\mathscr{A})$. Also $A \in m(\mathscr{A})$ because
	$A_n \uparrow A$. Hence $A \in M_0$.
	The argument for $A_n \downarrow A$ is identical (then $A_n^c \uparrow A^c$).

	Thus $M_0$ is a monotone class containing $\mathscr{A}$, so by minimality
	$m(\mathscr{A}) \subseteq M_0$. In particular, $m(\mathscr{A})$ is closed under complements.

	(iii) (Closure under countable unions.) First we show closure under finite unions.

	Define
	\[
		M_1 := \{A \in m(\mathscr{A}) : \forall\,B \in \mathscr{A},\; A \cup B \in m(\mathscr{A})\}.
	\]
	Clearly $M_1 \subseteq m(\mathscr{A})$. We check that $M_1$ is a monotone class containing
	$\mathscr{A}$.

	\emph{(a) $\mathscr{A} \subseteq M_1$.}
	If $A \in \mathscr{A}$ and $B \in \mathscr{A}$, then $A \cup B \in \mathscr{A} \subseteq m(\mathscr{A})$
	(since $\mathscr{A}$ is a field). Hence $A \in M_1$, so $\mathscr{A} \subseteq M_1$.

	\emph{(b) $M_1$ is a monotone class.}
	Suppose $A_n \uparrow A$ with $A_n \in M_1$ and fix a $B \in \mathscr{A}$.
	Then $A_n \cup B \uparrow A \cup B$. Since each $A_n \cup B \in m(\mathscr{A})$ and $m(\mathscr{A})$ is
	a monotone class, $A \cup B \in m(\mathscr{A})$. $B$ was arbitrary in $\mathscr{A}$, so $A \in M_1$.
	The decreasing case is the same.

	Therefore $M_1$ is a monotone class containing $\mathscr{A}$, hence $m(\mathscr{A}) \subseteq M_1$.
	Equivalently,
	\[
		\forall\,A \in m(\mathscr{A}),\ \forall\,B \in \mathscr{A},\qquad A \cup B \in m(\mathscr{A}).
	\]

	Now define
	\[
		M_2 := \{A \in m(\mathscr{A}) : \forall\,B \in m(\mathscr{A}),\; A \cup B \in m(\mathscr{A})\}.
	\]
	Again $M_2 \subseteq m(\mathscr{A})$. We show $M_2$ is a monotone class containing $\mathscr{A}$.

	\emph{(a) $\mathscr{A} \subseteq M_2$.}
	Fix $A \in \mathscr{A}$ and $B \in m(\mathscr{A})$. From the conclusion about $M_1$ (with roles swapped),
	$B \cup A \in m(\mathscr{A})$, hence $A \cup B \in m(\mathscr{A})$. Thus $A \in M_2$.

	\emph{(b) $M_2$ is a monotone class.}
	If $A_n \uparrow A$ with $A_n \in M_2$ and $B \in m(\mathscr{A})$, then
	$A_n \cup B \uparrow A \cup B$. Since each $A_n \cup B \in m(\mathscr{A})$ and $m(\mathscr{A})$
	is a monotone class, $A \cup B \in m(\mathscr{A})$, so $A \in M_2$.
	The decreasing case is identical.

	Thus $M_2$ is a monotone class containing $\mathscr{A}$, so $m(\mathscr{A}) \subseteq M_2$.
	In particular, $m(\mathscr{A})$ is closed under finite unions.

	Finally, let $(A_n)_{n\ge1} \subseteq m(\mathscr{A})$ and set $U_m := \bigcup_{n=1}^m A_n$.
	By finite-union closure, $U_m \in m(\mathscr{A})$ for each $m$, and $U_m \uparrow \bigcup_{n=1}^\infty A_n$.
	Since $m(\mathscr{A})$ is a monotone class, $\bigcup_{n=1}^\infty A_n \in m(\mathscr{A})$.

	We have shown $m(\mathscr{A})$ is a $\sigma$-field containing $\mathscr{A}$, hence
	$\sigma(\mathscr{A}) \subseteq m(\mathscr{A})$. Combined with the first inclusion,
	$m(\mathscr{A}) = \sigma(\mathscr{A})$.
\end{proof}
\boxednote{This and the previous theorem give us two different ways to generate a $\sigma$-field. We can start with something more complicated (field) or less complicated (a $\pi$-system). In this course, the monotone class theorem will be more commonly used. The choice comes down to convenience.}
\begin{theorem}
	($\pi$-$\lambda$.) If $\mathscr{P}$ is a $\pi$-system, then \[
		l(\mathscr{P}) = \sigma(\mathscr{P})
	.\] 
\end{theorem}

\section{$\sigma$-fields}


\begin{definition}[Topology]
	Let $\Omega$ be a set. A \emph{topology} on $\Omega$ is a collection $\mathcal{T}\subseteq 2^\Omega$ such that:
	\begin{enumerate}[label=(\roman*)]
		\item $\emptyset\in\mathcal{T}$ and $\Omega\in\mathcal{T}$,
		\item if $U,V\in\mathcal{T}$, then $U\cap V\in\mathcal{T}$,
		\item if $\{U_i\}_{i\in I}\subseteq\mathcal{T}$, then $\bigcup_{i\in I}U_i\in\mathcal{T}$.
	\end{enumerate}
	The sets in $\mathcal{T}$ are called \emph{open sets}, and $(\Omega,\mathcal{T})$ is a \emph{topological space}.
\end{definition}

\begin{definition}[Measurable space]
	A \emph{measurable space}\boxednote{A measure space is a measurable space equipped with a measure.} is a pair $(\Omega,\mathscr{F})$ where $\Omega$ is a set and $\mathscr{F}\subseteq 2^\Omega$ is a $\sigma$-algebra on $\Omega$, i.e.:
	\begin{enumerate}[label=(\roman*)]
		\item $\Omega\in\mathscr{F}$,
		\item if $A\in\mathscr{F}$ then $A^c\in\mathscr{F}$,
		\item if $A_1,A_2,\ldots\in\mathscr{F}$ then $\bigcup_{n=1}^\infty A_n\in\mathscr{F}$.
	\end{enumerate}
\end{definition}

\begin{definition}[Borel $\sigma$-algebra]
	If $(\Omega,\mathcal{T})$ is a topological space, the \emph{Borel $\sigma$-algebra} on $\Omega$ is
	\[
	\mathscr{B}(\Omega) \coloneqq \sigma(\mathcal{T}),
	\]
	the smallest $\sigma$-algebra containing all open sets.
	In particular, for $\R$ with its usual topology,
	\[
	\mathscr{B}_{\R} \coloneqq \mathscr{B}(\R)=\sigma\big(\{U\subseteq \R : U \text{ is open}\}\big)
	=\sigma\big(\{[a,\infty): a \in \R\}\big).
	\]
\end{definition}

\begin{definition}[Measure]
	Suppose $\mathscr{E}$ is a set collection such that $\O \in \mathscr{E} \subseteq  2^{\Omega}$. We say $\mu: \mathscr{E}\to [0,\infty]$ is a measure if:
	\begin{enumerate}[label=(\roman*)]
		\item $\mu(\O) = 0$ 
		\item For disjoint $A_n \in \mathscr{E}$ such that $\cup_n A_n \in \mathscr{E}$, \[
			\mu(\cup_n A_n) = \sum_n \mu(A_n)
		.\] 
	\end{enumerate}
	In particular, if $\mu(\Omega) < \infty$, $\mu$ is called a finite measure. If $\mu(\Omega) = 1$, then it is a probability measure. If for any event $A \in \mathscr{E}$, there exists $A_n \in \mathscr{E}$ such that $\cup_n A_n \supseteq A$ and $\mu(A_n) < \infty$ for each $n$, we call $\mu$ is a $\sigma$-finite measure. Any set in a $\sigma$-finite measure is the countable union of sets each of finite measure.
\end{definition}

\begin{definition}[Probability space]
	We call $(\Omega, \mathscr{F}, \mu)$ a probability space if $\mu(\Omega) = 1$.
\end{definition}

\begin{remark}[Properties of measure]
	Suppose $\mu$ is a measure on a $\sigma$-field $\mathscr{E}$.\footnote{Note now $\mu$ is not a premeasure on an arbitrary set collection; it is defined on a $\sigma$-field for these properties.} Then
	\begin{enumerate}[label=(\roman*)]
		\item (Monotonicity.) If $A \subseteq B$, then $\mu(A) \le \mu(B)$.
		\item (Continuity from below.) If $(A_n)$ is an increasing sequence ($A_n \subseteq A_{n+1}$), then
		\[
			\lim_{n \to \infty} \mu(A_n) = \mu\!\left( \lim_{n \to \infty} A_n \right)
			\qquad\text{where }\ \lim_{n\to\infty}A_n:=\bigcup_{n=1}^\infty A_n .
		\]
		\item (Continuity from above.) If $(A_n)$ is a decreasing sequence ($A_{n+1} \subseteq A_n$) and $\mu(A_1) < \infty$, then
		\[
			\mu\!\left( \lim_{n \to \infty} A_n \right)= \lim_{n \to \infty} \mu(A_n)
			\qquad\text{where }\ \lim_{n\to\infty}A_n:=\bigcap_{n=1}^\infty A_n .
		\]
	\end{enumerate}
\end{remark}

\begin{proof}
	\begin{enumerate}[label=(\roman*)]
		\item If $A\subseteq B$, then $B = A \amalg (B \setminus A)$, so
		\[
			\mu(B) = \mu(A) + \mu(B \setminus A) \ge \mu(A)
		\]
		by nonnegativity.
		
		\item Assume $A_n \uparrow$ and set $A:=\bigcup_{n=1}^\infty A_n$. Define
		\[
			D_1 := A_1,\qquad D_m := A_m \setminus A_{m-1}\ \ (m\ge 2),
		\]
		(with the convention $A_0:=\varnothing$). Then the $D_m$ are pairwise disjoint and for each $n$,
		\[
			A_n = \amalg_{m=1}^n D_m,
			\qquad\text{and}\qquad
			A=\amalg_{m=1}^\infty D_m .
		\]
		By countable additivity,
		\[
			\mu(A)=\sum_{m=1}^\infty \mu(D_m).
		\]
		On the other hand, for each $n$,
		\[
			\mu(A_n)=\mu\!\left(\amalg_{m=1}^n D_m\right)=\sum_{m=1}^n \mu(D_m),
		\]
		so taking $n\to\infty$ gives
		\[
			\lim_{n\to\infty}\mu(A_n)=\lim_{n\to\infty}\sum_{m=1}^n \mu(D_m)=\sum_{m=1}^\infty \mu(D_m)=\mu(A)
			= \mu\!\left(\bigcup_{n=1}^\infty A_n\right).
		\]
		
		\item Assume $A_n \downarrow$ and $\mu(A_1)<\infty$. Consider $B_n := A_1 \setminus A_n$.
		Then $(B_n)$ is increasing, so by (ii),
		\[
			\mu\!\left(\bigcup_{n=1}^\infty B_n\right)=\lim_{n\to\infty}\mu(B_n).
		\]
		But
		\[
			\bigcup_{n=1}^\infty B_n
			= \bigcup_{n=1}^\infty (A_1\setminus A_n)
			= A_1 \cap \bigcup_{n=1}^\infty A_n^c
			= A_1 \cap \left(\bigcap_{n=1}^\infty A_n\right)^c
			= A_1 \setminus \bigcap_{n=1}^\infty A_n .
		\]
		Also, since $A_n\subseteq A_1$ and $\mu(A_1)<\infty$,
		\[
			\mu(B_n)=\mu(A_1\setminus A_n)=\mu(A_1)-\mu(A_n).
		\]
		Therefore,
		\[
			\mu\!\left(A_1 \setminus \bigcap_{n=1}^\infty A_n\right)
			= \lim_{n\to\infty}\bigl(\mu(A_1)-\mu(A_n)\bigr)
			= \mu(A_1)-\lim_{n\to\infty}\mu(A_n).
		\]
		Finally, since
		\[
			\mu\!\left(A_1 \setminus \bigcap_{n=1}^\infty A_n\right)=\mu(A_1)-\mu\!\left(\bigcap_{n=1}^\infty A_n\right),
		\]
		again using $\mu(A_1)<\infty$, we conclude
		\[
			\mu\!\left(\bigcap_{n=1}^\infty A_n\right)=\lim_{n\to\infty}\mu(A_n).
		\]
	\end{enumerate}
\end{proof}

\section{Constructing the Uniform Measure}


Given a premeasure\footnote{For example, a finitely/countably additive set function like one involving volumes of rectangles.} on a small family of sets (a finite algebra/field), how can we get an extension to a useful $\sigma$-field?

We will first construct a measure $\mu$ on a field $\mathscr{A}$.
Then we will extend $\mu$ from $\mathscr{A}$ to $2^{\Omega}$ through outer measure as an approximation of $\mu$.
Finally, we will restrict the domain to $\mathscr{F}_\tau \subseteq 2^{\Omega}$ so that $\mathscr{F}_\tau$ is a $\sigma$-field and $\mu$ is a measure on that field.

Ultimately, we want to construct the uniform measure on $\left( [0,1], \mathscr{B}_{[0,1]} \right)$.

\begin{example}
	With $\Omega = [0,1]$, put $\mathscr{A} = \{\text{finite unions of intervals}\}$. It is easy to construct a measure on $\mathscr{A}$ by considering lengths of disjoint intervals in its elements.
\end{example}

How do we extend the measure from $\mathscr{A}$ to $\sigma(\mathscr{A})$? If we are able to do so, it is unique? We first answer the second question.

\begin{theorem}
	Suppose $\mathscr{A}$ is a field and $\mu$ and $\nu$ are two finite measures on $\sigma(\mathscr{A})$ such that for every $A \in \mathscr{A}$, $\mu(A) = \nu(A)$. Then for every $B \in \sigma(\mathscr{A})$, \[
		\mu(B) = \nu(B)
	.\] 
\end{theorem}

\begin{proof}
	Define \[
		\mathscr{H} = \{A \in \sigma(\mathscr{A}): \mu(A) = \nu(A)\}
	.\] We know that $\mathscr{A} \subseteq  \mathscr{H} \subseteq \sigma(\mathscr{A})$.
	Suppose we can show that $\mathscr{H}$ is a monotone class. Then it must contain the monotone class generated by $\mathscr{A}$, which is $\sigma(\mathscr{A})$ by the monotone class theorem. This would show equality. Let's prove it!

	Suppose we have a monotone sequence $A_n$ in $\mathscr{H}$. Then by definition of $\mathscr{H}$,\boxednote{Note the existence of $\lim\limits_{n \to \infty} \mu(A_n), \nu(A_n)$ is guaranteed in $[0,\infty]$ by the sequence being monotone.}
	\begin{align*}
		\mu(A_n) &= \nu(A_n) \\
		\implies \lim\limits_{n \to \infty} \mu(A_n) &= \lim\limits_{n \to \infty} \nu(A_n) \\
		\mu( \lim\limits_{n \to \infty} A_n) &= \nu \left( \lim\limits_{n \to \infty} A_n \right) \\
		\implies \lim\limits_{n \to \infty} A_n &\in \mathscr{H}
	,\end{align*}
	so $\mathscr{H}$ is a monotone class.
\end{proof}

\begin{definition}[Outer measure]
	We\boxednote{The definition of outer measure differs from that of measure in three ways: (i) outer measure is defined on $2^{\Omega}$ instead of a subset; (ii) outer measure must explicitly be monotone, while measure does not; (iii) outer measure is only required to be countably \emph{sub}additive.} say that $\tau: 2^{\Omega} \to [0,\infty]$ is an outer measure if:
	\begin{enumerate}[label=(\roman*)]
		\item $\tau(\O) = 0$,
		\item If $A \subseteq B$, then $\tau(A) \le \tau(B)$, and
		\item $\tau\left( \cup_n A_n \right) \le \sum_n \tau(A_n)$.
	\end{enumerate}
\end{definition}

\begin{theorem}[Construction of outer measure]\boxednote{Think of $\mathscr{E}$ as a field, even though the theorem does not require this. For instance, $\mathscr{E}$ might be all finite unions of rectangles.} 
	Suppose $\O \in \mathscr{E} \subseteq 2^{\Omega}$ and $\mu: \mathscr{E}\to [0,\infty]$ is such that $\mu(\O) = 0$. For every $A \in 2^{\Omega}$, define \[
		\tau(A) = \inf \{\sum_n \mu(A_n) : A_n \in \mathscr{E},\ A \subseteq \cup_n A_n\}
	.\] 
	\boxednote{If there does not exist a countable covering of $A$, convention is to put $\tau(A) = \infty$.}Then $\tau$ is an outer measure.
\end{theorem}

\begin{proof}
	\begin{enumerate}[label=(\roman*)]
		\item Since $m(\O) = 0$ and $\O \subseteq \O \cup \O \cup \ldots$, \[
			\tau(\O) \le \sum\limits_{n=1}^{\infty}  m(\O) = 0
			.\] $\tau$ is an $\inf$ over nonnegative numbers, so $\tau(\O) = 0$.
		\item Any covering of $B$ is also a covering of $A$. 
			\item For any $\epsilon > 0$ and $A_n$, there exists $\{B_{n,k}\}_{k}$ such that $B_{n,k} \in \mathscr{E}$, $\cup_k B_{n,k} \supseteq A_n$, and
			\[
				\tau(A_n) + \frac{\epsilon}{2^{n}} \ge \sum_k \mu(B_{n,k}).
			\]

		We know \[
			\bigcup\limits_{n}^{} A_n \subseteq \bigcup\limits_{n}^{} \bigcup\limits_{k}^{} B_{n,k}
		,\] so \[
			\tau\left( \bigcup\limits_{n}^{} A_n \right) \le \sum\limits_{n}^{} \sum\limits_{k}^{} \mu(B_{n,k}) \le \sum\limits_{n}^{} \tau(A_n) + \epsilon 
		.\] We can send $\epsilon \to 0$ to conclude.
	\end{enumerate}

\end{proof}

Now that we have an outer measure on $2^{\Omega}$, how do we restrict the domain to where we have countable additivity of disjoint sets?

\begin{theorem}[Restriction to Carath\'eodory-measurable sets]
	Let $\tau$ be an outer measure. Define \[
		\mathscr{F}_\tau := \{A \subseteq \Omega: \forall\ B \subseteq  \Omega, \tau(B) = \tau(B \cap A) + \tau(B \cap A^{c})\}
	.\]

	Then $\left( \Omega, \mathscr{F}_\tau, \tau \right)$ is a measure space.
\end{theorem}

Thinking of $A$ as a knife that splits any set $B$ into $A\cap B$ and $A^{c} \cap B$, we may think that $A$ is a good knife if it divides all sets in $\Omega$ well.

\begin{proof}
	Step I ($\mathscr{F}_\tau$ is a $\sigma$-field). We see that $\O \in \mathscr{F}_\tau$, as for any $B \supseteq \Omega$, \[
		\tau(B) = \underbrace{\tau(B \cap \O)}_{= \tau(\O) = 0} + \underbrace{\tau(B \cap \O^{c})}_{=\tau(B)}
	.\] Furthermore, because if $\tau(B) = \tau(A \cap B) + \tau(A^{c} \cap B) = \tau(A^{c} \cap B) + \tau(A \cap B)$, $\mathscr{F}_\tau$ is closed under complementation.

	We will now show that if $A_1,A_2 \in \mathscr{F}_\tau$, $A_1 \cap A_2 \in \mathscr{F}_\tau$.\boxednote{There is no good reason for proving closure under intersection instead of union. The proof for showing closure under finite union is analogous, so it might be a good idea to do it as an exercise.} For any $B \in \Omega$, 
	\begin{align*}
		\tau(B) &= \tau(B \cap A_1) + \tau(B \cap A_1^{c}) \\
				&= \tau(B\cap A_1 \cap A_2) + \underbrace{\tau(B\cap A_1 \cap A_2^{c}) + \tau(B \cap A_1^{c} \cap A_2) + \tau(B \cap A_1^{c} \cap A_2^{c})}_{\text{goal: make into } B \cap \left( A_1\cap A_2 \right)^{c}}
	.\end{align*}
	Showing the union of these three terms is equal to $B \cap \left( A_1 \cap A_2 \right)^{c}$ is a bit difficult, so let's try decomposing the single set. 
	\begin{align*}
		\tau\left(B \cap \left( A_1 \cap A_2 \right)^{c}\right) &= \tau \left( B \cap \left( A_1^{c} \cup A_2^{c} \right) \right) \\
																&= \underbrace{\tau \left( B \cap (A_2^{c} \cap A_1 )\right)}_{\text{one of goal terms}} + \tau \left( B \cap \left( A_1^{c} \right) \right) \\
		&= \tau\left(B \cap \left( A_2^{c} \cap A_1 \right)\right) + \tau (B \cap A_1^{c} \cap A_2) + \tau (B \cap A_1^{c} \cap A_2^{c})
	,\end{align*}
	where we cut with $A_1$ for the second equality and by $A_2$ for the third.

	Thus $\mathscr{F}_\tau$ is closed under finite intersection. Because we already showed it is closed under complementation, it is closed under finite union as well.

	It just remains to show that $\mathscr{F}_\tau$ is closed under countable union. Suppose $A_n \in \mathscr{F}_\tau$. WLOG, we can assume that $\{A_n\}$ is pairwise disjoint. We know from $\mathscr{F}_\tau$ being closed under finite union that for any $m$, $\cup_{n=1}^{m} A_n \in \mathscr{F}_\tau$. For any $B \subseteq \Omega$, 
	\begin{align*}
		\tau(B) &= \tau \left( B \cap \left( \bigcup\limits_{n=1}^{m} A_n \right) \right) + \tau \left( B \cap \left( \bigcup\limits_{n=1}^{m} A_n \right)^{c} \right)\\
		&= \tau(B\cap A_1) + \tau\left( B \cap \left( \bigcup\limits_{n=2}^{m} A_n \right) \right) + \tau \left( B \cap \left( \bigcup\limits_{n=1}^{m} A_n  \right)^{c} \right)
	,\end{align*}
	since we can cut the first term by $A_1$ and the $A_n$s are disjoint. Iterate on this. Then by monotonicity of outer measure,
	\begin{align*}
		\tau(B) &= \sum\limits_{n=1}^{m} \tau(B \cap A_n) + \tau\left( B \cap \left( \bigcup\limits_{n=1}^{m} A_n \right)^{c} \right) \\
				&\ge \sum\limits_{n=1}^{m} \tau(B \cap A_n) + \tau \left( B \cap \underbrace{\left( \bigcup\limits_{n=1}^{\infty} A_n \right)^{c}}_{\text{now smaller}} \right) \\
				&\ge \sum\limits_{n=1}^{\infty} \tau(B \cap A_n) + \tau \left( B \cap \left( \bigcup\limits_{n=1}^{\infty} A_n \right)^{c} \right)	\\
				&\ge \tau\left(B \cap \left( \bigcup\limits_{n=1}^{\infty} A_n \right)\right) + \tau \left( B \cap \left( \bigcup\limits_{n=1}^{\infty} A_n \right)^{c}  \right)
	,\end{align*}
	where the last inequality follows from countable subadditivity of $\tau$. However, also by countable (in this case, finite) subadditivity, \[
		\tau(B) \le \tau\left(B \cap \left( \bigcup\limits_{n=1}^{\infty} A_n \right)^{c}\right) + \tau \left( B \cap \left( \bigcup\limits_{n=1}^{\infty} A_n \right)^{c} \right)
	,\] 
	so $\cup_{n=1}^{\infty}A_n$ cuts arbitrary $B \subseteq \Omega$ well, so $\cup_{n=1}^{\infty}A_n \in \mathscr{F}_\tau$ and $\mathscr{F}_\tau$ is closed under countable union.

	Step II ($\tau$ is a measure on $\mathscr{F}_\tau$). We check $\tau(\O) = 0$ because it is an outer measure.	The only other condition for $\tau$ being a measure is countable additivity on disjoint sets.

	Again by $\tau$ as an outer measure, for disjoint $\{A_n\}$, \[
		\tau\left( \cup_n A_n \right) \le \sum_n \tau(A_n)
	.\] For the other direction, take $B = \cup_n A_n$, so
	\begin{align*}
		\tau(B) = \tau(\cup_n A_n) &\le \sum\limits_{n=1}^{\infty} \tau(A_n) \\
								   &= \sum\limits_{n=1}^{\infty} \tau(A_n \cap B)  + \tau\left( A_n \cap B^{c} \right) \\
								   &\le \sum\limits_{n=1}^{\infty} \tau(A_n) + \tau(0) \\
								   &= \sum\limits_{n=1}^{\infty} \tau(A_n)
	.\end{align*}

	$\mathscr{F}_\tau$ being a $\sigma$-field on $\Omega$ and $\tau$ being a valid measure on $\mathscr{F}_\tau$ satisfies the definition of $(\Omega, \mathscr{F}_\tau, \tau)$ being a measure space, so we are done.
\end{proof}

\begin{theorem}[Unique extension of measure]
	Suppose $\mu$ is a finite measure on a field $\mathscr{A}$. Then there exists a unique extension of $\mu$ from $\mathscr{A}$ to $\sigma(\mathscr{A})$.
\end{theorem}

\begin{proof}
	For every $A \subseteq \Omega$, define \[
		\tau(A) = \inf \{\sum_n \mu(B_n): B_n \in \mathscr{A},\ \cup_n B_n \supseteq A\}
	.\] By our earlier theorem, $\tau$ is an outer measure.

	Defining $\mathscr{F}_\tau$ according to the Caratheodory criterion, we know $(\Omega, \mathscr{F}_\tau, \tau)$ is a measure space.

	Since we already showed that a measure on a field is unique on a $\sigma$-algebra generated by that field, we just need to show that $\tau$ is actually an extension of $\mu$ to $\sigma(\mathscr{A})$\boxednote{This is why we need $\mu$ to be a measure on a \emph{field}.}. This entails showing that for any $A \in \mathscr{A}$, $\mu(A) = \tau(A)$ and that $\mathscr{F}_\tau \supseteq  \sigma(\mathscr{A})$.

	Fix $A \in \mathscr{A}$. Since $A \subseteq A \cup \O \cup \O \cup \ldots$, \[
		\tau(A) \le \mu(A) + \O + \O + \cdots = \mu\left(A \right)
	.\] For the other direction of the inequality, note that for any covering $B_n$ of $A$, 
	\begin{align*}
		\mu(A) &= \mu\left( \cup_n \left( A \cap B_n \right) \right) \\
			   &\le \sum_n \mu(A \cap B_n) \\
			   &\le \sum_n \mu(B_n)
	,\end{align*}
	where the second inequality follows from monotonicity of $\mu$.
	Taking an $\inf$ yields \[
		\mu(A) \le \tau(A)
	.\] 

	Since $\mathscr{F}_\tau$ is a $\sigma$-field, in order to show that $\mathscr{F}_\tau \supseteq \sigma(\mathscr{A})$, we just need to prove that $A \subseteq \mathscr{F}_\tau$. Hence, we want to show for every $A \in \mathscr{A}$ and every $B \subseteq \Omega$, $\tau(B) = \tau(B \cap A) + \tau(B \cap A^{c})$. The $\le$ direction is given by the subadditivity of $\tau$ with $\left( B \cap A \right) \cup \left( B \cap A^{c} \right)$

	Fix $\epsilon > 0$. There exists $B_n \in \mathscr{A}$ such that $\cup_n B_n \supseteq B$ and \[
		\tau(B) + \epsilon \ge \sum_n \mu(B_n)
	.\] 
		Then\boxednote{The last two inequalities use only outer-measure properties. At this point we do not yet know, for example, that $B \cap A \in \mathscr{F}_\tau$.}
	\begin{align*}
		\tau(B) + \epsilon &\ge \sum\limits_{n}^{} \mu(B_n) \\
				&= \sum_n \left[ \mu(B_n \cap A) + \mu(B_n \cap A^{c}) \right] \text{ by additivity of measure} \\
						   &= \sum_n \left[ \tau(B_n \cap A) + \tau(B_n \cap A^{c}) \right]  \text{ since $B_n, A, A^{c} \in \mathscr{A}$, where $\mu,\tau$ agree} \\
						   &\ge \tau\left( \left( \cup_n B_n \right)\cap A \right) + \tau\left( \left( \cup_n B_n\right)\cap A^{c} \right) \text{ by subadditivity of outer measure} \\
						   &\ge \tau\left( B\cap A \right) + \tau(B \cap A^{c}) \text{ by monotonicity of outer measure}
	.\end{align*}
	Sending $\epsilon \to 0$, \[
		\tau(B) \ge \tau(B \cap A) + \tau\left( B \cap A^{c} \right)
	,\] so we have shown \[
		\tau(B) = \tau(B \cap A) + \tau(B \cap A^{c})
	\] for arbitrary $A \in \mathscr{A}$ and $B \subseteq \Omega$, so $\mathscr{A} \subseteq \mathscr{F}_\tau$, and hence $\sigma(\mathscr{A}) \subseteq \mathscr{F}_\tau$.
\end{proof}

Now we finally consider $\left( [0,1], \mathscr{B}_{[0,1]} \right)$ and $\mathscr{A}$ as the set of finite unions of intervals and $\mu$ the length of an interval. By the theorem, we can extend $\mu$ to be a measure on $\mathscr{F} \supseteq \sigma(\mathscr{A}) = \mathscr{B}_{[0,1]}$.

\begin{definition}[Complete measure space]
	A measure space $(\Omega, \mathscr{F},\mu)$ is complete if \[
		A \in \mathscr{F} \text{ and } \mu(A) = 0 \implies \forall\ B \subseteq A, B \in \mathscr{F},
	.\] i.e. if the $\sigma$-field contains all subsets of measure zero sets.
\end{definition}

\begin{theorem}
	Given an outer measure $\tau$, the measure space $(\Omega, \mathscr{F}_\tau, \tau)$ constructed from the Caratheodory criterion is a complete measure space.
\end{theorem}

\begin{proof}
	We already showed that $(\Omega, \mathscr{F}_\tau, \tau)$ is a measure space. Suppose $\tau(A) = 0$. For any $B \in \Omega$, \[
		\tau(B) \le \tau(B \cap A) + \tau(B \cap A^{c})
	\] is true by subadditivity.

		By monotonicity of outer measure,
		\[
			\tau(B \cap A^{c}) \le \tau(B)
		.\]
		Also, $\tau(B \cap A) \le \tau(A) = 0$, so
		\[
			\tau(B) \ge \tau(B \cap A) + \tau(B \cap A^{c})
		.\] 

	Thus $\tau(B) = \tau(B \cap A) + \tau(B \cap A^{c})$, for every $B \in \Omega$, so $A \in \mathscr{F}_\tau$.
\end{proof}

We receive the fact that $\left| \mathscr{B}_{[0,1]} \right| = \left| [0,1] \right|$, which allows us to claim that $\left( [0,1], \mathscr{B}_{[0,1]}, \mu \right)$, where $\mu$ is the uniform measure, is not a complete measure space.\footnote{This implies that the measure space we constructed via Caratheodory is strictly larger than this Borel measure space.}

\begin{definition}[Cantor set]
	Define $C_0 = [0,1]$, $C_1 = [0,\frac{1}{3}] \cup [\frac{2}{3},1]$, and recursively $C_n = \frac{C_{n-1}}{3} \cup \left( \frac{2}{3} + \frac{C_{n-1}}{3} \right)$. Then define the Cantor set $C = \cap_{n=0}^{\infty}C_n$. 
\end{definition}

Immediately, $C \in \mathscr{B}_{[0,1]}$, and $\mu(C) = \lim\limits_{n \to \infty} \left( \frac{2}{3} \right)^{n} = 0$. However, $\left| C \right| = \left| [0,1] \right|$ because $C$ bijects with binary expansions of numbers in $[0,1]$. Thus\boxednote{The inequality is because for any set $S$, $\left| 2^{S} \right| > \left| S \right|$.} \[\left| 2^C \right| = \left| 2^{[0,1]} \right| > \left| [0,1] \right| = \left| \mathscr{B}_{[0,1]} \right|.\] Because this inequality is strict, there exists $A \in 2^C$ such that $A \not\in \mathscr{B}_{[0,1]}$.

What is the connection between the result $\mathscr{F}_\tau$ of Caratheodory extension and $\mathscr{B}_{[0,1]}$?

\begin{theorem}[Completion]
	Let $\left( \Omega, \mathscr{F},\mu \right)$ be a measure space. Define \[
		\tilde{\mathscr{F}} = \{A \cup N : A \in \mathscr{F} \text{ and } \exists\  B \in \mathscr{F} \text{ s.t. } \mu(B) = 0 \text{ and } B \supseteq N\}
	.\] 
	$\mathscr{F}$ is a $\sigma$-field. For any \boxednote{$\tilde{\mathscr{F}}$ just adds all subsets of sets in $\mathscr{F}$ with measure $0$.}$A \cup N \in \tilde{\mathscr{F}}$, defining \[
		\tilde{\mu}(A \cup N ) = \mu(A)
	,\] $\left( \Omega, \tilde{\mathscr{F}}, \tilde{\mu} \right)$ is a complete and well-defined measure space. Moreover, for any $A \in \mathscr{F}$, $\tilde{\mu}(A) = \mu(A)$; the complete measure space is an extension of the original measure space.

	We call $\tilde{\mathscr{F}}$ the \emph{completion} of $\mathscr{F}$.\boxednote{$\tilde{\mathscr{F}}$ is the smallest completion of $\mathscr{F}$.}
\end{theorem}

\begin{proof}
	We need to check that $\tilde{\mathscr{F}}$ is a $\sigma$-field.
	\begin{enumerate}[label=(\roman*)]
		\item Since $\O = \O \cup \O$, where $\O \in \mathscr{F}$ and $\O \subseteq \O$, $\O \in \tilde{\mathscr{F}}$
		\item Suppose $A \in \mathscr{F}$ and $N \subseteq B \in \mathscr{F}$ such that $m(B)=  0$.\footnote{With this sentence, we are supposing $A \cup N \in \tilde{\mathscr{F}}$.} We need to show $\left( A \cup N \right)^{c} \in \tilde{F}$. WE want to write $\left( A \cup N \right)^{c}$ as the union of something in $\mathscr{F}$ and a subset of a set of measure $0$.

			We can write\footnote{draw} \[
				\left( A \cup N \right)^{c} = \left( A \cup B \right)^{c} \cup \left( B \setminus \left( A \cup N \right) \right)
			.\] Since $A,B \in \mathscr{F} \implies \left( A \cup B \right)^{c} \in \mathscr{F}$ and $B \setminus \left( A \cup N \right) \subseteq B$ with $\mu(B) = 0$, we are done.
		\item Suppose $A_n \in \mathscr{F}$, $N_n \subseteq B_n \in \mathscr{F}$, and $\mu(B_n) = 0$. We want to show $\cup_n \left( A_n \cup N_n \right) \in \tilde{\mathscr{F}}$.

			Writing \[
				\cup_n \left( A_n \cup N_n \right) = \left( \cup_n A_n \right)\cup \left( \cup_n N_n \right)
			,\] since $\cup_n N_n \subseteq \cup_n B_n$ and $\mu(\cup_n B_n) = 0$, we are done.
	\end{enumerate}

	We want to show $\tilde{\mu}$ is well-defined. Suppose $A_1,A_2 \in \mathscr{F}$, $N_1 \subseteq B_1 \in \mathscr{F}$, $N_2 \subseteq B_2 \in \mathscr{F}$, $\mu(B_1) = \mu(B_2) = 0$, and $A_1 \cup N_1 = A_2 \cup N_2$. Can we show $\tilde{\mu}(A_1 \cup N_1) = \tilde{\mu}(A_2 \cup N_2)$?

	By definition,\boxednote{Adding in $B_1$ comes from: (i) monotonicity giving $\mu(A_1) \le \mu(A_1 \cup B_1)$ and (ii) subadditivity giving the other direction.}
	\begin{align*}
		\tilde{\mu}(A_1 \cup B_1) &= \mu(A_1) \\
		&= \mu(A_1 \cup B_1) \text{ because we can union with measure zero set} \\
		&\ge \mu(A_2) = \tilde{\mu}(A_2 \cup N_2)
	.\end{align*} By symmetry, the other direction is true, so $\tilde{\mu}$ is well-defined.

	Now we show $\tilde{\mu}$ is a measure on $\tilde{\mathscr{F}}$. 
	\begin{enumerate}[label=(\roman*)]
		\item $\tilde{\mu}(\O) = \tilde{\mu}\left( \O \cup \O \right) = \mu(\O) = 0$
		\item Suppose $A_n \in \mathscr{F}$, $N_n \subseteq B_n \in \mathscr{F}$, $\mu(B_n) = 0$, and $A_n \cup N_n$ are pairwise disjoint.

			Then 
			\begin{align*}
				\tilde{\mu}\left( \cup_n \left( A_n \cup N_n \right) \right) &= \tilde{\mu}\left( \left( \cup_n A_n \right)\cup \left( \cup_n N_{n} \right) \right) \\
																			 &= \mu\left( \cup_n A_n \right) \\
																			 &= \sum_n \mu(A_n) \text{ by disjointness} \\
																			 &= \sum_n \tilde{\mu}(A_n \cup N_n)
			,\end{align*}
	\end{enumerate}
	so $\tilde{\mu}$ is a measure.

	Since $\tilde{\mu}(A) = \tilde{\mu}\left( A \cup \O \right) = \mu(A)$, $\tilde{\mu}$ is an extension of $\mu$.
\end{proof}

\begin{theorem}
	Suppose\boxednote{One example is taking $\mathscr{A}$ as all finite unions of intervals in the unit interval and $\mathscr{F}_\tau$ as the \emph{completion} of the Borel $\sigma$-field.} $\mu$ is a finite measure on a field $\mathscr{A}$ of subsets of $\Omega$ and that $\tau$ is the outer measure constructed by $\mu$. With $\mathscr{F}_\tau$ as the set of $\tau$-measurable subsets of $\Omega$ according to the Caratheodory criterion, $(\Omega, \mathscr{F}_\tau, \tau)$ is the completion of $\left( \Omega, \sigma(\mathscr{A}), \tau \right)$.
\end{theorem}

We require the following lemma.

\begin{lemma}
	Under the same setting (finite measure $\mu$ with induced outer measure $\tau$), for any $A \in \mathscr{F}_\tau$, there exists $B \in \sigma(\mathscr{A})$ such that $B \supseteq A$ and $\tau(B) = \tau(A)$. Elements in $\mathscr{F}_\tau$ can be approximated above by elements in $\sigma(\mathscr{A})$.
\end{lemma}

\begin{proof}
	By definition of outer measure induced by a premeasure, \[
		\tau(A) = \inf \{\sum_n \mu(B_n) : B_n \in \mathscr{A}, \cup_n B_n \subseteq A\}
	.\] For every $m \in \N$, there exists a sequence $\{B_{m,n}\}_n \subseteq \mathscr{A}$ such that $\cup_n B_{m,n} \supseteq A$ and \[
		\tau(A) + \frac{1}{n} \ge \sum_n \mu(B_{m,n})
	.\] 

	Write \[
		B = \cap_m \cup_n B_{m,n} \supseteq A
	.\] Then $B \in \sigma(\mathscr{A})$ because each $B_{m,n} \in \mathscr{A}$.
	We see
	\begin{align*}
		\tau(A) + \frac{1}{m} &\ge \sum_n \mu(B_{m,n}) \\
							  &= \sum_n \tau(B_{m,n}) \\
							  & \ge \tau \left( \cup_n B_{m,n} \right) \\
							  &\ge \tau(B)			
	.\end{align*}
	Sending $m \to \infty$, $\tau(A) \ge \tau(B)$. Since $A \subseteq B$, $\tau(A) \le \tau(B)$.
\end{proof}

\begin{proof}
	(Of the theorem.) Write $\mathscr{F} = \sigma(\mathscr{A})$. We want to show that $\tilde{\mathscr{F}} = \mathscr{F}_\tau$, where \[
		\tilde{\mathscr{F}} = \{A \cup N : A \in \mathscr{F}, \exists\  B \in \mathscr{F} \text{ s.t. } B \supseteq  N, \mu(B) = 0\}
	.\] We know $\tilde{\mathscr{F}} \subseteq \mathscr{F}_\tau$ since $\tilde{\mathscr{F}}$ is the smallest completion of $\mathscr{F}$ and $\mathscr{F}_\tau$ is a completion of $\mathscr{F}$.

	For $\mathscr{F}_\tau \subseteq \tilde{\mathscr{F}}$, fix $A \in \mathscr{F}_\tau$. We want to write $A$ as the union of something in $\mathscr{F}$ and a subset of a measure zero set.

	By the lemma, there exists $B \in \mathscr{F}$ such that $B \supseteq A$ and $\tau(B) = \tau(A)$. By finiteness of $\tau$, this implies $\tau(B \setminus A) = 0$.

	Apply the lemma again to $B \setminus A$: there exists $C \in \mathscr{F}$ such that $C \subseteq B \setminus A$ and $\tau(C) = \tau(B \setminus A) = 0$.

	Draw, then write \[
		A = \left( A \cap C^{c} \right) \cup \left( A\cap C \right)
	.\] We don't know that $A \cap C^{c} \in \mathscr{F}$, because we just know $A \in \mathscr{F}_\tau$. However, we can write \[
		A = (B \cap C^{c}) \cup (A \cap C)
	.\] $B, C \in \mathscr{F}$, so $B \cap C^{c} \in \mathscr{F}$. Furthermore, $A \cap C \subseteq C$, where $\tau(C) = 0$, so we are done.
\end{proof}
